\section{Detailed lower-bound proof}
\label{sec:lower-proof}

\subsection{Electrostatic estimate with exact point nuclei}
First separate the Coulomb operator before full Wick ordering:
\begin{equation}
H_C^{\mathrm{dens}}=
\frac{\alpha}{2}D_L(\rho,\rho)+\alpha D_L(\rho,\nu)
+\alpha\sum_{i<j}Z_iZ_jG_L(R_i-R_j).
\label{eq:density-coulomb}
\end{equation}
The first term is positive, but the other terms require a uniform estimate.

Fix \(s>0\), eventually \(s=\Lambda_f^{-2}\), and write
\begin{equation}
G_{L,s}=e^{s\Delta}G_L,\qquad
h_{L,s}=G_L-G_{L,s},\qquad \nu_s=e^{s\Delta}\nu.
\end{equation}
If \(p_{L,t}\) denotes the torus heat kernel, then
\begin{equation}
h_{L,s}(x)=4\pi\int_0^s\bigl(p_{L,t}(x)-V^{-1}\bigr)\dd t.
\end{equation}
Positivity and unit integral of the heat kernel imply
\begin{equation}
h_{L,s}(x)\ge-\frac{4\pi s}{V},\qquad
\|h_{L,s}\|_{L^1}\le8\pi s.
\label{eq:heat-bounds}
\end{equation}
A lattice-shell estimate gives, uniformly in \(L\),
\begin{equation}
G_{L,2s}(0)=\frac1V\sum_{q\ne0}
\frac{4\pi e^{-2s|q|^2}}{|q|^2}\le Cs^{-1/2}.
\label{eq:heat-self}
\end{equation}
After writing \(q=2\pi n/L\), this reduces to
\(\sum_{n\ne0}|n|^{-2}e^{-a|n|^2}\le C a^{-1/2}\).
Counting lattice points in successive shells proves this estimate.

An exact rearrangement of \eqref{eq:density-coulomb} is
\begin{align}
H_C^{\mathrm{dens}}={}&
\frac{\alpha}{2}D_L(\rho+\nu_s,\rho+\nu_s)
\nonumber\\
&+\alpha D_L(\rho,\nu-\nu_s)
\nonumber\\
&+\alpha\sum_{i<j}Z_iZ_jh_{L,2s}(R_i-R_j)
\nonumber\\
&-\frac{\alpha}{2}G_{L,2s}(0)\sum_iZ_i^2.
\label{eq:heat-identity}
\end{align}
The first line is positive. Heat smoothing here is solely an identity used
in the proof; it does not replace the point-source Hamiltonian.
By \eqref{eq:local-bounds} and \eqref{eq:heat-bounds},
\begin{equation}
\left|D_L(\rho,\nu-\nu_s)\right|
\le16\pi d_Ls Z_{\mathrm{tot}}.
\label{eq:cross-heat}
\end{equation}
The absolute value denotes the corresponding operator bound. The third
line of \eqref{eq:heat-identity} is bounded below by
\begin{equation}
-\frac{4\pi\alpha s}{V}Z_{\mathrm{tot}}^2.
\end{equation}
Admissibility \eqref{eq:capacity} gives
\begin{equation}
\frac{Z_{\mathrm{tot}}^2}{V}\le2d_LZ_{\mathrm{tot}}.
\end{equation}
Finally, \(\sum_iZ_i^2\le Z_{\max}Z_{\mathrm{tot}}\le Z_{\max}^2K\).
Combining these estimates,
\begin{equation}
H_C^{\mathrm{dens}}\ge
-C\alpha\left(d_LsZ_{\max}+s^{-1/2}Z_{\max}^2\right)K.
\end{equation}
For \(L\ge\Lambda_f^{-1}\) and \(s=\Lambda_f^{-2}\),
\begin{equation}
H_C^{\mathrm{dens}}\ge-A_\Lambda K,\qquad
A_\Lambda\le C\alpha\Lambda_f(Z_{\max}+Z_{\max}^2).
\label{eq:electrostatic-bound}
\end{equation}
This is an actual uniform electrostatic estimate, including arbitrarily
close distinct nuclei.

\subsection{Cost of full Wick ordering}
We next account for full Wick ordering. This cannot simply be ignored.
In the finite sea representation define
\begin{equation}
n(A)=c^\dagger Ac-\Tr(P_-A).
\end{equation}
For one-particle matrices \(A,B\), Wick expansion gives
\begin{equation}
\begin{split}
n(A)n(B)={}&W(A,B)+n(M_{A,B})\\
&+\Tr\bigl(P_-A(1-P_-)B\bigr),
\end{split}
\label{eq:wick-identity}
\end{equation}
where \(W(A,B)\) denotes the fully Wick-ordered quartic expression and
\begin{equation}
M_{A,B}=AB-AP_-B-BP_-A.
\end{equation}
Equation \eqref{eq:wick-identity} follows directly by contracting the four
CAR operators.

For the compressed Fourier multipliers
\(T_q=P_fe^{-iqx}P_f\otimes I_{\C^4}\), one has \(\|T_q\|\le1\).
Thus the difference in \eqref{eq:wick-identity} has norm at most \(CM_L\).
For example, \(\|M_{A,B}\|\le3\) when \(\|A\|,\|B\|\le1\),
\(\|n(M_{A,B})\|\le8M_L\|M_{A,B}\|\), and the remaining trace has
absolute value at most \(4M_L\).
Moreover,
\begin{equation}
\frac1V\sum_{0<|q|\le2\Lambda_f}\frac{4\pi}{|q|^2}
\le C\Lambda_f.
\label{eq:coulomb-mode-sum}
\end{equation}
Therefore
\begin{equation}
\|H_C-H_C^{\mathrm{dens}}\|
\le C\alpha M_L\Lambda_f\le B_\Lambda V,\qquad
B_\Lambda\le C\alpha\Lambda_f^4,
\end{equation}
and
\begin{equation}
H_C\ge-A_\Lambda K-B_\Lambda V.
\label{eq:full-coulomb-bound}
\end{equation}
This controls the complete Coulomb interaction over all allowed pair
sectors. The volume term is a fixed consequence of the chosen ordering,
not a state-dependent subtraction.

\subsection{Transverse field and the infrared loss}
For the transverse field, \eqref{eq:square} gives
\begin{equation}
H_\gamma+H_I\ge
-\frac{e^2}{2V}\sum_{k,\lambda}
\frac{(\varepsilon_\lambda(k)\cdot j_k)^\dagger
(\varepsilon_\lambda(k)\cdot j_k)}{|k|^2}.
\label{eq:photon-current}
\end{equation}
Suppose every coupled photon mode satisfies \(|k|\ge k_{\min}\).
Polarization projection and Parseval yield
\begin{align}
H_\gamma+H_I
&\ge-\frac{e^2}{2k_{\min}^2}
\int_{\T_L^3}\sum_i j_i(x)^2\dd x
\nonumber\\
&\ge-\frac{6e^2d_L^2V}{k_{\min}^2}.
\label{eq:photon-bound}
\end{align}
The Parseval identity is an operator quadratic-form identity for a finite
Fourier polynomial; it does not require commuting current components.

For ordinary torus photons, \(k_{\min}=2\pi/L\). Since \(H_D\ge0\),
\eqref{eq:full-coulomb-bound}--\eqref{eq:photon-bound} prove
\begin{equation}
H_{\Lambda,L}\ge-A_\Lambda K-B_\Lambda V-D_\Lambda VL^2,
\qquad D_\Lambda\le C\alpha\Lambda_f^6.
\label{eq:lower-final}
\end{equation}
For an additional fixed infrared cutoff \(k_{\min}\ge\kappa>0\), they instead
prove
\begin{equation}
H_{\Lambda,L}^{(\kappa)}
\ge-A_\Lambda K-
\left(B_\Lambda+C\alpha\Lambda_f^6\kappa^{-2}\right)V.
\label{eq:ir-final}
\end{equation}
This proof is nonperturbative. It controls arbitrary photon occupation and
every admissible pair sector. Its failure to be extensive without \(\kappa\)
is not merely a weakness of the estimate.

\section{Counterexamples, failed comparisons, and stress tests}
\label{sec:counterexamples}

\subsection{A neutral-current counterexample at large volume}
Set \(K=0\). Work in the equivalent finite bare-Dirac CAR representation.
The free sea \(P_-\) has rank \(2M_L\).
Let
\begin{equation}
q=\frac{2\pi}{L},\qquad C_L=P_f\cos(qx_1)P_f
\end{equation}
on the scalar cutoff space. Define \(S_L=\sign(C_L)\), taking
\(\sign(0)=+1\). Then \(S_L\) is a self-adjoint unitary.
Define the one-particle projector
\begin{equation}
\gamma_L=\frac12\left(I+\mathsf a_2S_L\right).
\label{eq:trial-projector}
\end{equation}
The spatial and spin operators commute, so \(\gamma_L^2=\gamma_L\).
Since \(\tr_{\C^4}\mathsf a_2=0\), \(\rank\gamma_L=2M_L\).
Its Slater determinant therefore belongs to the \emph{exact neutral charge
sector}, not merely a state of zero expected charge.
Its mean charge density vanishes pointwise:
\begin{equation}
\tr_{\C^4}\gamma_L(x,x)=2M_L/V
=\tr_{\C^4}P_-(x,x).
\label{eq:trial-charge}
\end{equation}
Its transverse current, however, is macroscopic:
\begin{align}
J_L&:=\int_{\T_L^3}\cos(qx_1)\langle j_2(x)\rangle\dd x
\nonumber\\
&=\Tr\bigl((\gamma_L-P_-)\mathsf a_2C_L\bigr)
=2\Tr|C_L|.
\label{eq:trial-current}
\end{align}
The free-sea contribution vanishes because \(C_L\) has zero momentum
diagonal. Since \(\|C_L\|\le1\),
\begin{align}
J_L&\ge2\Tr C_L^2
\nonumber\\
&=\#\{p\in\Gamma_{f,L}:p+qe_1\in\Gamma_{f,L}\}
\nonumber\\
&\ge c_{\Lambda_f}V
\label{eq:trial-current-size}
\end{align}
for sufficiently large \(L\). The last inequality follows by counting
points away from the boundary of the cutoff ball.

The free fermionic energy obeys
\begin{equation}
\langle H_D\rangle\le4M_L\sqrt{\Lambda_f^2+m^2}\le C_\Lambda V.
\label{eq:trial-free}
\end{equation}
The Coulomb energy is also only \(O(V)\), including charge fluctuations.
To see this directly for full Wick ordering, put \(Q_L=\gamma_L-P_-\).
Quasifree Wick expansion gives
\begin{equation}
\begin{split}
\langle:\rho_{-k}\rho_k:\rangle
={}&\Tr(Q_LT_{-k})\Tr(Q_LT_k)\\
&-\Tr(Q_LT_{-k}Q_LT_k).
\end{split}
\label{eq:trial-wick}
\end{equation}
The first term vanishes by \eqref{eq:trial-charge}. Also,
\begin{equation}
\left|\Tr(Q_LT_{-k}Q_LT_k)\right|
\le\|Q_L\|_{\HS}^2\le4M_L.
\end{equation}
Using \eqref{eq:coulomb-mode-sum},
\begin{equation}
|\langle H_C\rangle|\le C_\Lambda V.
\label{eq:trial-coulomb}
\end{equation}
Now take a product with a photon coherent state whose mean vector potential is
\begin{equation}
A(x)=a\,e_2\cos(qx_1).
\end{equation}
This field is transverse. Its modes are retained once \(q\le\Lambda_\gamma\).
Its photon energy and current coupling are
\begin{equation}
\frac14Vq^2a^2+eaJ_L.
\end{equation}
Minimizing in \(a\) gives
\begin{equation}
-\frac{e^2J_L^2}{Vq^2}\le-c_\Lambda e^2VL^2.
\label{eq:trial-gain}
\end{equation}
Combining \eqref{eq:trial-free}, \eqref{eq:trial-coulomb}, and
\eqref{eq:trial-gain} proves
\begin{equation}
\boxed{E^{\vac}_{\Lambda,L}\le-c_\Lambda L^5+C_\Lambda L^3.}
\label{eq:vacuum-instability}
\end{equation}
Here \(c_\Lambda>0\) for any fixed nonzero coupling.

Every trial state is a legitimate finite-volume Fock-space state. Its
photon occupation can be large, but no photon-number restriction was
imposed. Total charge is exactly zero, and Coulomb fluctuations have been
included.
The optimized vector potential is \(O(L^2)\), while its magnetic field is
\(O(L)\). A momentum cutoff alone does not exclude these large fields.
Thus calling the model low-energy QED cannot remove this counterexample
unless an additional field or energy restriction is explicitly imposed.
Fixed mass counterterms contribute only \(O(V)\) to these states. A fixed
vacuum-energy-density counterterm also contributes \(O(V)\). A fixed
nonzero charge renormalization changes the coefficient in
\eqref{eq:vacuum-instability}, not its \(L^5\) behavior.

\subsection{What the construction proves and what it does not}
Equation \eqref{eq:vacuum-instability} disproves \eqref{eq:W} in the stated
normalization. It also prevents an ordinary finite vacuum-energy density
for this regulator. It does not prove \(E(R,Z)-E^{\vac}<-CK\).
Subtracting the exact interacting vacuum energy removes the common
superextensive term. A comparison between the nuclear and vacuum ground
states is still needed.
That is the precise missing step in the primary proof:
\begin{quote}
Control the energy change caused by nuclei relative to the actual
interacting vacuum, uniformly in \(L\).
\end{quote}
Separate lower bounds on both systems do not provide that comparison.
Even bounds
\begin{equation}
E(R,Z)\ge-AK-BV,\qquad -BV\le E^{\vac}\le0
\end{equation}
would imply only \(E(R,Z)-E^{\vac}\ge-AK-BV\), not \(-CK\).

\subsection{A separate all-\texorpdfstring{\(L\)}{L} counterexample}
Under convention \eqref{eq:nuclear-self}, choose \(L\) small enough that
only fermion momentum \(p=0\) remains and no nonzero photon mode is retained.
Take two nuclei with \(K=2\) and \(Z_1=Z_2=1\).
Neutrality forces occupation of both electron spin states and no positrons.
All nonzero Fourier components of the fermion density vanish, so
\begin{equation}
E_{\Lambda,L}(R,Z)-E^{\vac}_{\Lambda,L}
=2m+\alpha G_L(R_1-R_2).
\label{eq:small-box}
\end{equation}
The zero-mean Green function \(G_1\) is positive near its singularity and
therefore negative somewhere else. Choose \(r_0\ne0\) with \(G_1(r_0)<0\).
Since \(G_L(Lr_0)=L^{-1}G_1(r_0)\), taking \(R_1-R_2=Lr_0\) gives
\begin{equation}
E_{\Lambda,L}-E^{\vac}_{\Lambda,L}\longrightarrow-\infty
\quad(L\downarrow0)
\end{equation}
at fixed \(K=2\).
Thus the literal theorem covering every box size is false for this
convention. This is a periodic small-box effect, not a macroscopic collapse
mechanism. A meaningful macroscopic formulation should impose
\(L\ge L_0>0\) and specify its finite self-energy convention.

\clearpage
\subsection{The requested stress tests}
\small
\noindent
\begin{tabular}{@{}>{\raggedright\arraybackslash}p{0.28\textwidth}>{\raggedright\arraybackslash}p{0.67\textwidth}@{}}
\toprule
Potential obstruction & Outcome in the stated model\\
\midrule
Point-Coulomb singularities &
Finite fermion projection makes the external operator bounded; nuclear
pair singularities remain positive. This provides no cutoff-independent
Coulomb estimate.\\[0.6em]
Collective nuclear concentration &
The heat-kernel estimate controls all distinct configurations uniformly
at fixed cutoff.\\[0.6em]
Incorrect negative-energy quantization &
Using \(d\Gamma(D_0)\) without filling and subtracting the sea is not our
positive electron--positron Hamiltonian; without a cutoff it is already
unbounded below.\\[0.6em]
Unlimited pair production &
At fixed cutoff, capacity is finite. Its growth as \(O(V)\) does not prevent
the \(O(-L^5)\) photon-current mechanism.\\[0.6em]
Vacuum polarization &
Included through the full Fock-space interactions, but no uniform estimate
relative to the interacting vacuum has been obtained.\\[0.6em]
Transverse magnetic interactions &
They produce the explicit obstruction \eqref{eq:vacuum-instability}.\\[0.6em]
UV and IR limits &
The constants diverge with \(\Lambda_f\); removing the auxiliary photon
infrared cutoff destroys estimate \eqref{eq:ir-final}.\\[0.6em]
Charge and Gauss constraints &
Exactly enforced through \eqref{eq:neutrality}--\eqref{eq:capacity};
neither counterexample uses a forbidden charge sector.\\[0.6em]
Finite-system versus extensive stability &
Every finite system is bounded below, while its lower spectral bound can
behave as \(-L^5\).\\
\bottomrule
\end{tabular}
\normalsize

\section{What the required literature actually establishes}
\label{sec:literature}
\subsection{Lieb--Yau (1988)}
Their model has fixed-number fermionic electrons with kinetic energy
\(\sqrt{p^2+m^2}-m\), point nuclei, and Coulomb interactions. It has no
positrons, pair creation, or photons. Their sufficient stability conditions
include
\begin{equation}
Z\alpha\le2/\pi,\qquad \alpha<1/(47q),
\end{equation}
where \(q\) is spin degeneracy. For \(q=2\), physical \(\alpha\) and \(Z\le26\)
satisfy these conditions. This is a rigorous stability theorem for a
different many-body Hamiltonian~\cite{LY1988}.

The sharp one-center obstruction in that model is transparent. If
\(\nu>2/\pi\), sharpness of the relativistic Coulomb inequality gives \(f\)
with
\begin{equation}
\langle f,(|p|-\nu/|x|)f\rangle<0.
\end{equation}
For \(f_t(x)=t^{3/2}f(tx)\),
\begin{equation}
\begin{split}
&\langle f_t,(\sqrt{p^2+m^2}-m-\nu/|x|)f_t\rangle\\
&\hspace{1em}=
t\langle f,(|p|-\nu/|x|)f\rangle+o(t),
\end{split}
\end{equation}
which tends to \(-\infty\). This scaling is excluded by a fixed momentum cutoff.

\subsection{Lieb--Loss (2002)}
Their model includes point nuclei and ultraviolet-cutoff quantized photons,
but is explicitly \emph{no-pair}. Electrons are restricted using the positive
spectral subspace of the field-dependent Dirac operator \(D(A)\). They prove
stability at \(\alpha=1/137\) for \(Z\le42\), so iron is included. The
field-dependent projection is substantive: it cannot be replaced by a fixed
free projection or by unrestricted electron--positron Fock space~\cite{LL2002}.

\subsection{Lieb (2002)}
This survey distinguishes stability of individual systems, stability of the
second kind, and thermodynamic questions. Its relativistic quantized-field
discussion concerns a positive-subspace model, not a proof for pair-producing
full QED~\cite{Lieb2002}.

\subsection{J\"akel--Wreszinski (2018)}
Their finite-torus cutoff model includes electrons, positrons, pair terms,
and quantized photons, but no external nuclei. Their stated bound is
\begin{equation}
H(\Lambda,V)+\mu(\Lambda,V)I\ge0
\end{equation}
for a finite volume-dependent \(\mu\). It does not establish \(\mu=O(V)\),
nor a nuclear energy bound relative to the interacting vacuum. Thus its
statement does not contradict \eqref{eq:vacuum-instability}. Our
finite-regulator lemma independently proves the semiboundedness needed here,
without importing formal photon-displacement identities~\cite{JW2018}.

\subsection{Hainzl--Lewin--S\'er\'e: polarized BDF vacua}
In a no-photon quasifree/mean-field setting, their variational estimates
include, for finite cutoff and \(0\le\alpha<4/\pi\),
\begin{equation}
\begin{split}
\mathcal E(Q)\ge{}&
(1-\alpha\pi/4)\Tr_{P_-}(D_0Q)\\
&+\frac{\alpha}{2}D(\rho_Q-n,\rho_Q-n)
-\frac{\alpha}{2}D(n,n).
\end{split}
\label{eq:BDF}
\end{equation}
They establish polarized-vacuum minimizers under the stated hypotheses.
The external density belongs to Coulomb space; a literal delta has infinite
Coulomb self-energy. At finite fermion cutoff the external coupling can be
represented using a filtered source, but nuclear repulsion and uniformity
still require separate work. Most importantly, a lower bound on a restricted
quasifree variational class is not a lower bound on arbitrary correlated
Fock states~\cite{HLS2004,HLS2005}.

\subsection{Point-nucleus Dirac--Coulomb theory}
For \(D_0-\nu/|x|\), essential self-adjointness on the usual punctured smooth
core holds through \(\nu=\sqrt3/2\). For \(\sqrt3/2<\nu<1\), a distinguished
extension exists with domain in \(H^{1/2}\). General multicenter
distinguished-extension results apply when each atomic charge weight is
below one. These operators still possess a negative continuum;
self-adjointness is not semiboundedness or matter stability. Distinct
iron-or-lighter nuclei are individually subcritical, while coalescence
combines their singular weights~\cite{ELS2021}.
These results address complementary pieces. None supplies the missing
full-Fock, massless-photon, vacuum-relative comparison in \eqref{eq:S}.

\section{Low-energy reduction: controlled steps and missing inputs}
\label{sec:effective}
For the free dispersion there is a useful rigorous estimate:
\begin{equation}
0\le\frac{p^2}{2m}-\bigl(\sqrt{p^2+m^2}-m\bigr)
\le\frac{p^4}{8m^3}.
\label{eq:dispersion}
\end{equation}
It follows from Taylor's theorem for \(\sqrt{1+x}\), or directly by
integrating its second derivative for \(x\ge0\).
On an \(N\)-electron subspace with every \(|p|\le p_*\),
\begin{equation}
\left\|\sum_{i=1}^N\left[
\frac{p_i^2}{2m}-\bigl(\sqrt{p_i^2+m^2}-m\bigr)
\right]\right\|
\le\frac{Np_*^4}{8m^3}.
\label{eq:free-error}
\end{equation}
This is an extensive free-particle error bound. It does not control pair
creation, strong fields, or radiative corrections.

For the interacting problem, let \(\Pi\) project onto a chosen no-positron
subspace and \(\Pi^\perp=1-\Pi\). Whenever the complementary resolvent
exists, the exact Feshbach operator is
\begin{equation}
\begin{split}
F_\Pi(H-z)={}&\Pi(H-z)\Pi\\
&-\Pi H\Pi^\perp
\left[\Pi^\perp(H-z)\Pi^\perp\right]^{-1}
\Pi^\perp H\Pi.
\end{split}
\label{eq:feshbach}
\end{equation}
If one could prove
\begin{equation}
\Pi^\perp(H-z)\Pi^\perp\ge g\Pi^\perp,\qquad g>0,
\label{eq:gap}
\end{equation}
then, on the appropriate quadratic-form domain,
\begin{equation}
\begin{split}
0\le{}&
\Pi H\Pi^\perp
\left[\Pi^\perp(H-z)\Pi^\perp\right]^{-1}
\Pi^\perp H\Pi\\
\le{}&g^{-1}\Pi H\Pi^\perp H\Pi.
\end{split}
\label{eq:feshbach-bound}
\end{equation}
This identifies two concrete missing inputs:
\begin{itemize}
\item a pair-sector gap \(g\) uniform in admissible nuclear configurations
and system size;
\item a bound on the virtual-pair correction that is extensive, rather than
an operator-norm estimate growing uncontrollably with volume.
\end{itemize}
A statement that energies lie below \(2m\) is insufficient. One must
specify the reference energy and prove the interacting threshold. Soft
photons remove a global spectral gap above the ground state, although they
do not by themselves eliminate the free pair-production threshold. A
separately established gap to states containing additional pairs could
still be useful.

After suitable block diagonalization and parameter matching, the expected
leading effective Hamiltonian is
\begin{equation}
H_{\mathrm{PF}}=
\sum_{i=1}^N
\frac{\left[\boldsymbol\sigma_i\cdot
(-i\nabla_i+e_{\mathrm{phys}}A(x_i))\right]^2}
{2m_{\mathrm{phys}}}
+V_C^{\mathrm{phys}}+H_\gamma,
\label{eq:PF}
\end{equation}
with relativistic and radiative corrections.
Equation \eqref{eq:PF} is a \emph{target effective description}, not a
consequence proved here. Controlling its remainder requires estimates on
covariant momenta, field strengths, and virtual-pair terms. A bound on
photon momentum alone does not bound field amplitude.

There are rigorous reductions in restricted models.
Hainzl--Lewin--S\'er\'e prove a nonrelativistic limit from BDF to
Hartree--Fock for fixed particle number and suitable external densities
and cutoff scaling. That is a no-photon mean-field result, not a
full-QED-to-Pauli--Fierz theorem with macroscopic error
control~\cite{HLS2008}. Arai's Dirac--Maxwell limit is another restricted
bridge: it concerns one Dirac particle coupled to photon Fock space,
rather than a quantized Dirac field with pair creation~\cite{Arai2003}.

To transfer effective-theory stability back to full QED, one would need both:
\begin{enumerate}
\item a lower comparison of the effective operator with a stable Coulomb
or Pauli--Fierz Hamiltonian, with error at most \(CK\);
\item control of the complementary pair sector, so that lower-energy
states have not been excluded by the reduction.
\end{enumerate}
A stable effective Hamiltonian alone provides neither requirement.

\section{Strongest justified partial conclusions and further requirements}
\label{sec:conclusions}
For the explicit massless, fixed-momentum-cutoff model, the strongest
global estimate established here is
\begin{equation}
\boxed{
E_{\Lambda,L}(R,Z)\ge
-A_\Lambda K-B_\Lambda L^3-D_\Lambda L^5,\qquad
L\ge\Lambda_f^{-1}.}
\end{equation}
The power \(L^5\) is necessary in free-vacuum normalization because of
the matching vacuum trial construction.
For the model with fixed infrared cutoff in the current coupling,
\begin{equation}
\boxed{
E_{\Lambda,L}^{(\kappa)}(R,Z)
\ge-A_\Lambda K-B_{\Lambda,\kappa}L^3.}
\end{equation}
This is a genuine extensive theorem with point nuclei, quantized photons,
and pair creation. It is stronger than finite-volume boundedness, but it
requires an additional regulator.

Since the free vacuum has expectation zero in the no-nucleus Hamiltonian,
\(E_{\Lambda,L}^{\vac,(\kappa)}\le0\). Consequently,
\begin{equation}
E_{\Lambda,L}^{(\kappa)}(R,Z)
-E_{\Lambda,L}^{\vac,(\kappa)}
\ge-A_\Lambda K-B_{\Lambda,\kappa}L^3.
\end{equation}
If one additionally imposes a positive lower nuclear density
\(K/V\ge\rho_*>0\), this implies
\begin{equation}
E^{(\kappa)}-E^{\vac,(\kappa)}
\ge-\left(A_\Lambda+B_{\Lambda,\kappa}/\rho_*\right)K.
\end{equation}
That conditional result does not cover arbitrarily dilute systems or
removal of \(\kappa\).

\subsection{Changing the ultraviolet cutoff at fixed physical parameters}
Our estimates grow roughly as \(\Lambda_f\), \(\Lambda_f^4\), and
\(\Lambda_f^6\). No cutoff-uniform conclusion follows. A nonperturbative
mass-and-charge matching procedure and a regulator preserving the
necessary gauge structure are additional requirements.

\subsection{Thermodynamic limit}
For the baseline model in free-vacuum normalization,
\begin{equation}
E^{\vac}_{\Lambda,L}/V\le-c_\Lambda L^2+C_\Lambda,
\end{equation}
so a finite vacuum-energy density fails. Subtracting the exact finite-volume
vacuum energy poses a different thermodynamic problem. Even the extensive
bound with \(\kappa>0\) does not by itself prove convergence; upper bounds,
screening, boundary control, and suitable subadditivity are also needed.

\subsection{Atoms and molecules}
Finite-box ground states follow from compact resolvent. Infinite-volume
bound states require binding inequalities and infrared compactness;
molecular binding additionally requires comparison with separated
clusters. None follows from a lower stability bound alone.

\subsection{Equilibrium crystals}
These require minimizing over nuclear configurations, proving binding and
a selected positive density, and establishing crystalline order and
mechanical stability. An energetic lower bound contains no such structural
information.

\subsection{Requirements for a physical proof}
For the intended physical theorem, a promising program must begin by
addressing two distinct issues. First, choose and analyze a regulator
that controls large transverse fields without the fixed-momentum-cutoff
pathology demonstrated above. Second, prove a localized comparison with
the \emph{interacting} vacuum that charges the energy change to the nuclei
rather than to the whole volume.
The required theorem would then have to control pair-sector spectral
thresholds, vacuum polarization, and low-energy reduction errors uniformly
in system size. Those are substantive missing estimates. The no-pair
and BDF theorems provide useful components, but combining their conclusions
does not constitute a proof for the full model requested here.
